\section{Methodology}
\label{sec:methodology}

\subsection{Synthetic Instance Generators}

We evaluate on three families of synthetic HM-style constraint systems, each
designed to isolate a distinct structural property. We use synthetic generators
to keep the source, conflict structure, and weighting scheme fully reproducible.

\paragraph{Dataset A: Single Planted Conflict.}
Each instance contains exactly one conflict pair: two constraints binding the
same fresh type variable $\alpha$ to incompatible base types,
$\alpha = \texttt{Int}$ and $\alpha = \texttt{Bool}$, both with weight~$1.0$.
The remaining $|C| - 2$ constraints are satisfiable fillers over independent
fresh variables, including base, list, and function types.
This yields a unique MUS of size~2 and two symmetric minimum-weight correction
subsets, each of weight~$1.0$.

\paragraph{Dataset B: Overlapping Conflicts.}
Each instance plants $k = 4$ constraints binding a shared variable
$x$ to mutually incompatible types $\{\texttt{Int}, \texttt{Bool},
\texttt{Char}, \texttt{String}\}$ with weights $1.0, 2.0, 3.0, 4.0$
respectively.
Every pair of these constraints forms a MUS, yielding $\binom{k}{2} = 6$
pairwise MUSes, all involving the same shared variable.
The satisfiable fillers are appended as in Dataset~A.
This family tests whether MUS overlap alone is enough to produce an
approximation gap.

\paragraph{Dataset C: HM-Shaped Constructor Mismatch.}
Each instance plants a conflict between a function-type shape and a list-type
shape on the same variable $f$: the constraints
$f = [\alpha] \to \beta$, $\alpha = \texttt{Int}$, $\beta = \texttt{Bool}$,
and $f = [\texttt{Char}]$ are inconsistent because $f$ cannot simultaneously
be a function type and a list type.
Constraint weights are derived from type-expression depth:
\[
w_i = 1 + \max(\mathrm{depth}(\tau_i^L), \mathrm{depth}(\tau_i^R)).
\]
Fillers are randomly generated type expressions up to depth~$3$.

\subsection{Worked Example: Shared-Variable Overlap}

Case Study~2 illustrates the intended structure of Dataset~B using the program
\texttt{let x = 5 in (x ++ "hello", x \&\& True)}.
The variable $x$ receives a type variable $\tau$, and its three uses impose
the incompatible constraints $\tau = \texttt{Int}$, $\tau = [\texttt{Char}]$,
and $\tau = \texttt{Bool}$.
The resulting MUSes overlap on the same variable $\tau$.
This example supports the interpretation that Dataset~B contains genuine MUS
overlap, even though the overlap is too regular to produce an approximation gap.

\subsection{Greedy Approximation Algorithm}

We implement a lazy MUS-guided greedy hitting-set algorithm.
At each iteration, the algorithm invokes a MUS finder on the current residual
constraint set to obtain one MUS~$M$, adds $M$ to the set of discovered MUSes,
and recomputes a weighted greedy hitting set over all MUSes discovered so far.
At each greedy hitting-set step, it selects the constraint
\[
c^* = \arg\min_c \frac{w_c}{|\{M' : c \in M'\}|},
\]
where the denominator counts currently unhit discovered MUSes containing $c$.
The selected hitting set is removed from the original constraint set, and the
process repeats until the residual constraint set is satisfiable.

MUS extraction uses an unsatisfiable-core procedure followed by deletion-based
shrinking: constraints are tentatively removed one at a time and kept removed
if the subset remains unsatisfiable.
Satisfiability at each step is checked by Robinson unification over the
type-equality constraint system.
The full-information weighted hitting-set greedy algorithm has the standard
$H_d$ approximation bound, where $d$ is the maximum set size.
Our implementation uses this greedy rule lazily on discovered MUSes and is
therefore evaluated empirically against the exact baseline.

\subsection{Exact Baseline}

We use Z3's optimisation interface as the exact baseline.
Each constraint $c_i$ is guarded by a Boolean variable $\mathit{keep}_i$.
If $\mathit{keep}_i$ is true, the corresponding type equality must hold.
The objective minimises the total weight of dropped constraints, namely
constraints for which $\mathit{keep}_i$ is false.
This returns an exact minimum-weight correction subset for the encoded instance.
A timeout of 30{,}000\,ms per instance is applied; all 300~instances completed
within budget.

\subsection{Brute-Force Sanity Check}

To validate the Z3 baseline, we run exhaustive brute-force search on small
instances with $|C| = 10$.
This sanity check covers 15~instances, from 5~seeds across the 3~datasets.
The brute-force solver enumerates all $2^{|C|}$ candidate correction subsets
and returns the minimum-weight one.
Z3 and brute force agree on optimal weight for all 15~instances.

\subsection{Experimental Grid}

The full experiment spans 3~datasets, 5~sizes
$|C| \in \{10, 20, 50, 100, 200\}$, and 20~seeds, giving 300~instances.
For each instance we record greedy weight, exact weight, approximation ratio,
greedy and exact wall-clock runtime, number of greedy iterations, number of
MUSes discovered, MUS sizes, overlap measure, and weight statistics.
Code and results are available at
\url{https://github.com/Provia/hm-mcs-approx}.
